\documentclass[10pt,a4paper]{amsart}
\usepackage{amsmath,amssymb,amsthm,graphicx}
% Retained minimal preamble for the Stacks Project chapter more-morphisms.tex (More on Morphisms).
% Verbatim excerpt of the chapter preamble: Xy-pic setup (native preamble lines 14-18) and the native
% theorem declarations with the mathematical macros (native preamble lines 154-217). Copied mechanically;
% nothing added, removed or redefined. The wrapper class and the amsmath/amssymb/amsthm/graphicx base
% packages are supplied by the assembler, so the class/externaldocument part of the native preamble is omitted.
% For commutative diagrams we use Xy-pic
\usepackage[all]{xy}

% We use 2cell for 2-commutative diagrams.
\xyoption{2cell}
% Theorem environments.
%
\theoremstyle{plain}
\newtheorem{theorem}[subsection]{Theorem}
\newtheorem{proposition}[subsection]{Proposition}
\newtheorem{lemma}[subsection]{Lemma}

\theoremstyle{definition}
\newtheorem{definition}[subsection]{Definition}
\newtheorem{example}[subsection]{Example}
\newtheorem{exercise}[subsection]{Exercise}
\newtheorem{situation}[subsection]{Situation}

\theoremstyle{remark}
\newtheorem{remark}[subsection]{Remark}
\newtheorem{remarks}[subsection]{Remarks}

\numberwithin{equation}{subsection}

% Macros
%
\def\lim{\mathop{\mathrm{lim}}\nolimits}
\def\colim{\mathop{\mathrm{colim}}\nolimits}
\def\Spec{\mathop{\mathrm{Spec}}}
\def\Hom{\mathop{\mathrm{Hom}}\nolimits}
\def\Ext{\mathop{\mathrm{Ext}}\nolimits}
\def\SheafHom{\mathop{\mathcal{H}\!\mathit{om}}\nolimits}
\def\SheafExt{\mathop{\mathcal{E}\!\mathit{xt}}\nolimits}
\def\Sch{\mathit{Sch}}
\def\Mor{\mathop{\mathrm{Mor}}\nolimits}
\def\Ob{\mathop{\mathrm{Ob}}\nolimits}
\def\Sh{\mathop{\mathit{Sh}}\nolimits}
\def\NL{\mathop{N\!L}\nolimits}
\def\CH{\mathop{\mathrm{CH}}\nolimits}
\def\proetale{{pro\text{-}\acute{e}tale}}
\def\etale{{\acute{e}tale}}
\def\QCoh{\mathit{QCoh}}
\def\Ker{\mathop{\mathrm{Ker}}}
\def\Im{\mathop{\mathrm{Im}}}
\def\Coker{\mathop{\mathrm{Coker}}}
\def\Coim{\mathop{\mathrm{Coim}}}

% Boxtimes
%
\DeclareMathSymbol{\boxtimes}{\mathbin}{AMSa}{"02}

%
% Macros for moduli stacks/spaces
%
\def\QCohstack{\mathcal{QC}\!\mathit{oh}}
\def\Cohstack{\mathcal{C}\!\mathit{oh}}
\def\Spacesstack{\mathcal{S}\!\mathit{paces}}
\def\Quotfunctor{\mathrm{Quot}}
\def\Hilbfunctor{\mathrm{Hilb}}
\def\Curvesstack{\mathcal{C}\!\mathit{urves}}
\def\Polarizedstack{\mathcal{P}\!\mathit{olarized}}
\def\Complexesstack{\mathcal{C}\!\mathit{omplexes}}
% \Pic is the operator that assigns to X its picard group, usage \Pic(X)
% \Picardstack_{X/B} denotes the Picard stack of X over B
% \Picardfunctor_{X/B} denotes the Picard functor of X over B
\def\Pic{\mathop{\mathrm{Pic}}\nolimits}
\def\Picardstack{\mathcal{P}\!\mathit{ic}}
\def\Picardfunctor{\mathrm{Pic}}
\def\Deformationcategory{\mathcal{D}\!\mathit{ef}}
\usepackage{cleveref}
\crefname{originalreference}{original reference}{original references}
\Crefname{originalreference}{Original reference}{Original references}
\makeatletter
\newlabel{coherent-lemma-induction-principle}{{08DR}{source}{}{sourcepointer-1}{}}
\newlabel{coherent-lemma-induction-principle@cref}{{[originalreference][0][]08DR}{[1][0][]source}{}{}{}}
\makeatother
\crefname{originallemma}{original lemma}{original lemmas}
\Crefname{originallemma}{Original lemma}{Original lemmas}
\makeatletter
\newlabel{lemma-jumps-w}{{0F3G}{source}{}{sourcepointer-2}{}}
\newlabel{lemma-jumps-w@cref}{{[originallemma][0][]0F3G}{[1][0][]source}{}{}{}}
\makeatother
\crefname{originallemma}{original lemma}{original lemmas}
\Crefname{originallemma}{Original lemma}{Original lemmas}
\makeatletter
\newlabel{lemma-quasi-finite-separated-pass-through-finite}{{05K0}{source}{}{sourcepointer-3}{}}
\newlabel{lemma-quasi-finite-separated-pass-through-finite@cref}{{[originallemma][0][]05K0}{[1][0][]source}{}{}{}}
\makeatother
\crefname{originallemma}{original lemma}{original lemmas}
\Crefname{originallemma}{Original lemma}{Original lemmas}
\makeatletter
\newlabel{lemma-weighting-base-change}{{0F3B}{source}{}{sourcepointer-4}{}}
\newlabel{lemma-weighting-base-change@cref}{{[originallemma][0][]0F3B}{[1][0][]source}{}{}{}}
\makeatother
\crefname{originallemma}{original lemma}{original lemmas}
\Crefname{originallemma}{Original lemma}{Original lemmas}
\makeatletter
\newlabel{lemma-weighting-check-after-etale-base-change}{{0F39}{source}{}{sourcepointer-5}{}}
\newlabel{lemma-weighting-check-after-etale-base-change@cref}{{[originallemma][0][]0F39}{[1][0][]source}{}{}{}}
\makeatother
\crefname{originalreference}{original reference}{original references}
\Crefname{originalreference}{Original reference}{Original references}
\makeatletter
\newlabel{morphisms-lemma-massage-finite}{{03HW}{source}{}{sourcepointer-6}{}}
\newlabel{morphisms-lemma-massage-finite@cref}{{[originalreference][0][]03HW}{[1][0][]source}{}{}{}}
\makeatother
\crefname{originalreference}{original reference}{original references}
\Crefname{originalreference}{Original reference}{Original references}
\makeatletter
\newlabel{morphisms-theorem-chevalley}{{054K}{source}{}{sourcepointer-7}{}}
\newlabel{morphisms-theorem-chevalley@cref}{{[originalreference][0][]054K}{[1][0][]source}{}{}{}}
\makeatother
\crefname{originalreference}{original reference}{original references}
\Crefname{originalreference}{Original reference}{Original references}
\makeatletter
\newlabel{properties-lemma-stratification-locally-finite-constructible}{{0F2M}{source}{}{sourcepointer-8}{}}
\newlabel{properties-lemma-stratification-locally-finite-constructible@cref}{{[originalreference][0][]0F2M}{[1][0][]source}{}{}{}}
\makeatother
\title{Level sets of the integral of a weighting are locally constructible}
\author{The Stacks Project Authors}
\date{stacks-project:more-morphisms.tex:local-snapshot-20261011; GNU FDL 1.2 or later}
\begin{document}
\maketitle
\noindent\textit{Editorial note.} Counted span is the author's own proof environment, source lines 22625-22669, immediately following the statement at 22618-22623.\par\medskip
\noindent\textit{Editorial note.} The source line 22629 contains the author's typo 'weighthing'; it is retained verbatim inside the counted proof.\par\medskip
\noindent\textit{Editorial note.} Four cross-chapter pointers (Coherent, Morphisms, Properties) plus four same-chapter weightings lemmas are recorded in assembly.external\_references.\par\medskip
\noindent\textit{Allowed standard prerequisites (editorial).}\begin{itemize}
\item Level sets of a weighting are locally constructible. chapter lemma lemma-jumps-w, cited at proof line 22662. Transfers constructibility from w to the sum function in the trivial-residue-extension case.
\item Base change of the integral and of weightings. chapter lemmas lemma-weighting-check-after-etale-base-change and lemma-weighting-base-change, cited at proof lines 22626-22629. Shows the function of interest commutes with arbitrary base change, which licenses the surjective finite presentation reduction.
\item Chevalley's theorem and massaging a finite morphism. Morphisms, theorem morphisms-theorem-chevalley and lemma morphisms-lemma-massage-finite, cited at proof lines 22656-22667. Makes images of constructible sets constructible and reduces to trivial residue field extensions.
\item Stratification by the finite value set of w. Properties, lemma properties-lemma-stratification-locally-finite-constructible, cited at proof lines 22664-22665. Gives that w takes only finitely many values, so the integral is a finite sum of constructible pieces.
\item Induction principle for quasi-compact quasi-separated schemes. Coherent, lemma coherent-lemma-induction-principle, cited in the proof. Supports the reduction of the ambient scheme to a manageable piece.
\end{itemize}\medskip
\ifdefined\hypertarget\hypertarget{sourcepointer-1}{}\fi
\noindent\textit{Source reference.} Original reference 08DR: focused-seeds/geometry/S9-morph-jumps-int-w/ref-tags.txt, source lines 4--4. Navigation pointer only: the counted proof cites this statement by \textbackslash{}ref as a standing input; the deduction the counted proof has to make is carried out inside the counted span, and no part of that deduction is delegated to this pointer.\par\medskip
\ifdefined\hypertarget\hypertarget{sourcepointer-2}{}\fi
\noindent\textit{Source reference.} Original lemma 0F3G: sources/stacks-more-morphisms/pinned-original-source.tex, source lines 22514--22514. Navigation pointer only: the counted proof cites this statement by \textbackslash{}ref as a standing input; the deduction the counted proof has to make is carried out inside the counted span, and no part of that deduction is delegated to this pointer.\par\medskip
\ifdefined\hypertarget\hypertarget{sourcepointer-3}{}\fi
\noindent\textit{Source reference.} Original lemma 05K0: sources/stacks-more-morphisms/pinned-original-source.tex, source lines 12457--12457. Navigation pointer only: the counted proof cites this statement by \textbackslash{}ref as a standing input; the deduction the counted proof has to make is carried out inside the counted span, and no part of that deduction is delegated to this pointer.\par\medskip
\ifdefined\hypertarget\hypertarget{sourcepointer-4}{}\fi
\noindent\textit{Source reference.} Original lemma 0F3B: sources/stacks-more-morphisms/pinned-original-source.tex, source lines 22173--22173. Navigation pointer only: the counted proof cites this statement by \textbackslash{}ref as a standing input; the deduction the counted proof has to make is carried out inside the counted span, and no part of that deduction is delegated to this pointer.\par\medskip
\ifdefined\hypertarget\hypertarget{sourcepointer-5}{}\fi
\noindent\textit{Source reference.} Original lemma 0F39: sources/stacks-more-morphisms/pinned-original-source.tex, source lines 22111--22111. Navigation pointer only: the counted proof cites this statement by \textbackslash{}ref as a standing input; the deduction the counted proof has to make is carried out inside the counted span, and no part of that deduction is delegated to this pointer.\par\medskip
\ifdefined\hypertarget\hypertarget{sourcepointer-6}{}\fi
\noindent\textit{Source reference.} Original reference 03HW: focused-seeds/geometry/S9-morph-jumps-int-w/ref-tags.txt, source lines 5--5. Navigation pointer only: the counted proof cites this statement by \textbackslash{}ref as a standing input; the deduction the counted proof has to make is carried out inside the counted span, and no part of that deduction is delegated to this pointer.\par\medskip
\ifdefined\hypertarget\hypertarget{sourcepointer-7}{}\fi
\noindent\textit{Source reference.} Original reference 054K: focused-seeds/geometry/S9-morph-jumps-int-w/ref-tags.txt, source lines 6--6. Navigation pointer only: the counted proof cites this statement by \textbackslash{}ref as a standing input; the deduction the counted proof has to make is carried out inside the counted span, and no part of that deduction is delegated to this pointer.\par\medskip
\ifdefined\hypertarget\hypertarget{sourcepointer-8}{}\fi
\noindent\textit{Source reference.} Original reference 0F2M: focused-seeds/geometry/S9-morph-jumps-int-w/ref-tags.txt, source lines 7--7. Navigation pointer only: the counted proof cites this statement by \textbackslash{}ref as a standing input; the deduction the counted proof has to make is carried out inside the counted span, and no part of that deduction is delegated to this pointer.\par\medskip
\begin{lemma}
\label{lemma-jumps-int-w}
Let $f : X \to Y$ be a locally quasi-finite morphism of finite
presentation. Let $w : X \to \mathbf{Z}$ be a weighting of $f$. Then
the level sets of the function $\int_f w$ are locally constructible in $Y$.
\end{lemma}

\begin{proof}
By Lemma \ref{lemma-weighting-check-after-etale-base-change}
formation of the function $\int_f w$ commutes with arbitrary
base change and by Lemma \ref{lemma-weighting-base-change}
after base change we still have a weighthing.
This means that if we can find $Y' \to Y$
surjective and of finite presentation, then it
suffices to prove the result after base change to $Y'$, see
Morphisms, Theorem \ref{morphisms-theorem-chevalley}.

\medskip\noindent
The question is local on $Y$ hence we may assume $Y$ is affine.
Then $X$ is quasi-compact and quasi-separated
(as $f$ is of finite presentation). Suppose that $X = U \cup V$
are quasi-compact open. Then we have
$$
\textstyle{\int}_f w =
\textstyle{\int}_{f|_U} w|_U +
\textstyle{\int}_{f|_V} w|_V -
\textstyle{\int}_{f|_{U \cap V}} w|_{U \cap V}
$$
Thus if we know the result for $w|_U$, $w|_V$, $w|_{U \cap V}$
then we know the result for $w$. By the induction principle
(Cohomology of Schemes, Lemma \ref{coherent-lemma-induction-principle})
it suffices to prove the lemma when $X$ is affine.

\medskip\noindent
Assume $X$ and $Y$ are affine. We may choose an open immersion
$X \to T$ where $T \to Y$ is finite, see
Lemma \ref{lemma-quasi-finite-separated-pass-through-finite}.
Because we may still base change with a suitable $Y' \to Y$
we can use Morphisms, Lemma \ref{morphisms-lemma-massage-finite}
to reduce to the case where
all residue field extensions induced by the morphism $T \to Y$
(and a foriori induced by $X \to Y$) are trivial.
In this situation $\int_f w$ is just taking the sums
of the values of $w$ in fibres. The level sets of $w$
are locally constructible in $X$ (Lemma \ref{lemma-jumps-w}).
The function $w$ only takes a finite number of values by
Properties, Lemma
\ref{properties-lemma-stratification-locally-finite-constructible}.
Hence we conclude by
Morphisms, Theorem \ref{morphisms-theorem-chevalley}
and some elementary arguments on sums of integers.
\end{proof}

\end{document}
